\documentclass{ximera}

\title{Powers of Five}
\author{Kenneth Berglund}

\begin{document}
\begin{abstract}
    Investigating powers of 5.
\end{abstract}
\maketitle

We'll investigate powers of 5.

\begin{enumerate}
    \item Solve \(5^x = 25\). 
    \item Solve \(5^x = 125\).
    \item Solve \(5^x = 625\).
    \item Now we'll introduce a mystery number \(a\). The
    definition of \(a\) is that \(5^a = 50\). That is, \(a\)
    is the solution of the equation \(5^x = 50\).
    \begin{enumerate}
        \item Is \(a\) a whole number? (Whole numbers are 0, 1, 2, 3, 4, ...)
        \item What two consecutive integers does \(a\) fall between? (Some 
        examples of consecutive integers are \(-2\), \(-1\); or 5, 6; or 101, 102.)
    \end{enumerate}
\end{enumerate}

\begin{instructorNotes}
This really should hammer home that logarithms evaluated at positive numbers
do, in fact, produce real numbers. It tries to give students a place to 
put them on the number line. 
\end{instructorNotes}

\end{document}
